% chapters/lec02-dafny.tex
% COL 862/8595 --- Lecture 02, 30 July 2026 --- Subodh Sharma, IIT Delhi

\lstdefinelanguage{Dafny}{
  morekeywords={method,function,ghost,returns,requires,ensures,invariant,
    decreases,var,while,if,then,else,assert,nat,int,bool,return,forall,exists},
  sensitive=true, morecomment=[l]{//}, morestring=[b]"
}

\chapter{Invariants, Conditions, and Automated Reasoning via Dafny}

\section{Program analysis vs.\ formal verification}
Given a program $P$, program analysis asks whether a property $\phi$ holds for
all behaviours of $P$, written $P \models \phi$ (e.g.\ ``no divide by zero'').
It runs statically (data-flow analysis, abstract interpretation, symbolic
execution, model checking) or dynamically (testing, profiling, runtime
monitoring, taint analysis).

\begin{ppkey}
Formal Verification $=$ Model $+$ Specification $+$ Proof: the semantics of the
program, the correctness properties, and a technique proving the former meets
the latter.
\end{ppkey}

On a withdrawal routine, analysis \emph{infers} facts (\ppinline{amount}
influences \ppinline{balance}; there are two paths; no pointer is dereferenced).
Verification instead starts from a committed obligation,
$\mathit{balance} \geq 0 \wedge \mathit{amount} \geq 0 \Rightarrow
\mathit{balance}' \geq 0$, and proves it.

\begin{ppnote}
Both are hard for the reason Rice's theorem gives: every non-trivial semantic
property is undecidable. The standing tension is precision $\leftrightarrow$
scalability --- infinitely many inputs, exponentially many paths, unbounded
loops, and state involving pointers, heaps, recursion, exceptions and threads.
\end{ppnote}

\section{Where LLMs help, and where they do not}
LLMs assist with specification synthesis (natural language into logic, e.g.\
``every accepted request eventually receives a response'' as
$G(\mathit{req} \rightarrow F\,\mathit{resp})$), semantic inference of
invariants and pre-/post-conditions, and auxiliary lemmas when a solver stalls
on a fact it could in principle derive. Four layers: property discovery,
abstraction discovery, proof search, tool orchestration.

\begin{pppitfall}
LLMs cannot check their own output and are neither sound nor complete. They
reason correctly about $\forall x \exists y\,.\,R(x,y)$ in isolation, then
silently interchange the quantifiers a step later. Generate candidates; validate
with a verifier. \emph{Trust but verify.}
\end{pppitfall}

\section{Pre- and postconditions}
A \ppterm{precondition} holds immediately before a block runs and is the
\emph{caller's} responsibility; a \ppterm{postcondition} holds immediately after
and is the \emph{callee's}. A Hoare triple $\{P\}\,C\,\{Q\}$ combines them, e.g.
$\{x>0\}\; x := x*2 \;\{x \bmod 2 = 0\}$. The \ppterm{strongest postcondition}
is the tightest bound on the output state: from $\{x>2\}$, doubling gives
$\{x>4\}$, which is stronger than the stated $\{x \bmod 2 = 0\}$. Validity means
$\mathrm{sp} \Rightarrow Q$.

\begin{ppprogram}{Nomad Bridge: an incomplete precondition (USD 190M lost)}{prog:nomad}
\begin{lstlisting}[language=C]
mapping(bytes32 => uint256) public confirmAt; // Stores confirmed roots

function acceptableRoot(bytes32 _root) public view returns (bool) {
    // PRECONDITION CHECK: Is this root non-zero AND confirmed?
    return confirmAt[_root] != 0;
}

function process(bytes memory _message) external returns (bool) {
    bytes32 root = proveAndProcess(_message);
    require(acceptableRoot(root), "Invalid root!");  // enforcer
    // --- Critical State Change & Execution: payout ---
}
\end{lstlisting}
\end{ppprogram}

An upgrade accidentally set \ppinline{confirmAt[0] = 1}. An attacker submitted a
fake message; the Merkle tree returned root $0$; the check passed and funds were
paid out. The precondition needed was
$\_root \neq 0 \wedge \mathit{confirmAt}[\_root] > 0$ --- the bug is in the
specification, not the implementation.

\section{Weakest precondition}
$\mathrm{wp}(C,Q)$ is the least restrictive initial condition guaranteeing that
$C$ terminates with $Q$. For $\{P\}\;x := x+5\;\{x>15\}$, substitution gives
$x+5>15$, so $\mathrm{WP} = \{x>10\}$, computed by backward reasoning; any sound
$P$ satisfies $P \Rightarrow \mathrm{WP}$.

\begin{table}[htbp]
  \centering\small
  \begin{tabularx}{\linewidth}{@{}lX@{}}
    \toprule
    \textbf{Construct} & $\mathrm{wp}(S,Q)$ \\
    \midrule
    \ppinline{skip}   & $Q$ \\
    \ppinline{abort}  & $\mathrm{false}$ \\
    \ppinline{x := E} & $Q[x \mapsto E]$ \\
    \ppinline{S1; S2} & $\mathrm{wp}(S_1, \mathrm{wp}(S_2,Q))$ \\
    \ppinline{if B then S1 else S2}
      & $(B \implies \mathrm{wp}(S_1,Q)) \wedge (\neg B \implies \mathrm{wp}(S_2,Q))$ \\
    \ppinline{while B do S}
      & $I \wedge \forall\text{vars}.\big((I \wedge B \implies \mathrm{wp}(S,I))
        \wedge (I \wedge \neg B \implies Q)\big)$ \emph{(axiomatic approx.)} \\
    \bottomrule
  \end{tabularx}
  \caption{Weakest-precondition rules. This is what Dafny constructs automatically.}
  \label{tab:wp}
\end{table}

\begin{ppnote}
The loop rule is only an approximation: the true wp of a loop is an uncomputable
fixpoint, so the rule is parameterised by an invariant $I$ that must be
supplied. Its conjuncts are exactly the obligations below.
\end{ppnote}

\section{Invariants}
A \ppterm{loop invariant} $I$ holds before the loop, at the start of each
iteration, and on exit, and must satisfy:

\begin{itemize}
  \item \textbf{Initialisation}: $P \Rightarrow I$;
  \item \textbf{Preservation}: $I \wedge B \wedge \mathit{Body} \Rightarrow I'$;
  \item \textbf{Exit}: $I \wedge \neg B \Rightarrow Q$.
\end{itemize}

\begin{pppitfall}
Too weak and the invariant survives preservation but fails exit --- true and
useless. Too strong and preservation fails. Discovery is the search for a
predicate that is both inductive and strong enough to entail $Q$.
\end{pppitfall}

\section{Dafny}
\ppinline{requires} is a precondition, \ppinline{ensures} a postcondition,
\ppinline{invariant} a loop invariant, \ppinline{decreases} a variant for
termination.

\begin{ppprogram}{Dafny cannot discharge this without an invariant}{prog:dafny-twice}
\begin{lstlisting}[language=Dafny]
method Twice(n: nat) returns (x: nat, y: nat)
  ensures x == n
  ensures y == 2 * n
{
  x := 0;
  y := 0;

  while x < n
  {
    x := x + 1;
    y := y + 2;
  }
}
\end{lstlisting}
\end{ppprogram}

\begin{ppexample}
Adding \ppinline{invariant x <= n} and \ppinline{invariant y == 2 * x} closes
the proof. Initialisation: $x=y=0$. Preservation: $y=2x$ gives $y+2 = 2(x+1)$.
Exit: $\neg(x<n) \wedge x \leq n$ yields $x=n$, hence $y=2n$. Note that the
invariant supplies the relation between the mutable variables that the guard
converts into a statement about $n$.
\end{ppexample}

\begin{ppprogram}{GCD with a ghost specification}{prog:dafny-gcd}
\begin{lstlisting}[language=Dafny]
// Ghost function to define GCD mathematically for specification
ghost function Gcd(a: nat, b: nat): nat
  decreases b
{
  if b == 0 then a else Gcd(b, a % b)
}

method ComputeGcd(a: nat, b: nat) returns (res: nat)
  requires a > 0 || b > 0
  ensures res == Gcd(a, b)
{
  var x := a;
  var y := b;

  while y != 0
  {
    var temp := y;
    y := x % y;
    x := temp;
  }

  return x;
}
\end{lstlisting}
\end{ppprogram}

\begin{ppnote}
The variant is $y$: since $0 \leq x \bmod y < y$, and $(\mathbb{N},<)$ is
well-founded, the loop terminates. The invariant is
\ppinline{Gcd(x, y) == Gcd(a, b)}, preserved by the identity
$\gcd(x,y) = \gcd(y, x \bmod y)$; at exit $y = 0$ so $x = \mathrm{Gcd}(a,b)$.
Whether Z3 unfolds the recursive ghost definition far enough is the point of the
exercise.
\end{ppnote}

\section{Under the hood}
Dafny compiles programs into \ppterm{verification conditions} and passes the
\emph{negation} of the desired property to Z3: \ppinline{UNSAT} means no
counterexample exists and the VC is valid, \ppinline{SAT} exhibits a
counterexample, \ppinline{Unknown} means Z3 could not settle it. A reported
failure says $P \Rightarrow \mathrm{wp}(C,Q)$ was not proved --- either it is
false, or the invariant was too weak.

\begin{pppitfall}
\ppinline{Unknown} is neither a proof nor a counterexample. Reading a timeout as
success is the verification analogue of reading a green test suite as
correctness.
\end{pppitfall}

\section{Exercises}
\begin{ppexercises}
  \ppexercise{\emph{(Submitted work.)} Study the weakest precondition calculus ---
  its properties (law of the excluded miracle, monotonicity) and its rules ---
  and submit a worked example. Then compute the wp of Table~\ref{tab:wp} by hand
  on Program~\ref{prog:dafny-gcd}; write, scan and submit.}

  \ppexercise{Weaken the invariant of Program~\ref{prog:dafny-twice} to
  \ppinline{y <= 2 * x} and identify which of initialisation, preservation or
  exit now fails.}

  \ppexercise{State the precondition Program~\ref{prog:nomad} should have
  enforced, and argue why \ppinline{confirmAt[\_root] != 0} is strictly weaker.}

  \ppexercise{Exhibit a trace separating $G(\mathit{req}) \rightarrow
  F(\mathit{resp})$ from $G(\mathit{req} \rightarrow F\,\mathit{resp})$.}
\end{ppexercises}
